\itemhead{3.2 例题 · 例题4（教材p.84）}{XJ-B1-C03-EX-P084-Y0177}
\label{XJ-B1-C03-EX-P084-Y0177}
{\footnotesize 湘教版教材p.84 / PDF第91页 / 3.2 例题 / 例题4}
\field{原题}{判断下列函数的奇偶性：
\subq{1}$f(x)=x^2+|x|$；
\subq{2}$g(x)=x+\dfrac kx$；
\subq{3}$h(x)=x^3\ (x\in[-2,5])$。}
\field{参考答案}{（1）偶函数；（2）奇函数，按原分式定义域为 $\mathbb R\setminus\{0\}$（即使 $k=0$ 也不能自动补0）；（3）非奇非偶，因为定义域 $[-2,5]$ 不对称。}
{\small \field{关联知识}{F08 奇偶性、对称与函数性质应用（教材帮p.113、114、115 / PDF122、123、124）；F02 定义域与原式限制（教材帮p.79、80、81 / PDF88、89、90）}
\field{方法与易错点}{检查定义域，再代入 $-x$；求另一侧解析式时同时写出 $-x$ 所在区间。 易错：只看公式不看定义域；用一对点证明奇偶；负半轴代换遗漏最外层负号。}
\field{答案依据}{例题原解位置：湘教版教材p.84及题库原图；独立编写可编辑参考解答。答案册此部分主要收练习、习题，不伪造例题答案册位置。}}
